% =============================================================================
% chapter1.tex — Project Context and Foundation
% Chapter 1 of the Codara PFE Report
% =============================================================================

\chapter{Project Context and Foundation}
\label{ch:context}

% ---------------------------------------------------------------------------
\section{Introduction}
\label{sec:ch1-intro}
% ---------------------------------------------------------------------------

Every large-scale software migration project begins with the same uncomfortable
truth: the code that must be moved is precisely the code nobody fully
understands any more. In data-intensive enterprises, that code is most often
written in \ac{SAS} --- a proprietary, fourth-generation language that has
accumulated three decades of business logic, macro meta-programming patterns,
and procedural idioms that resist straightforward automated translation.
Rewriting these assets manually is slow, expensive, and error-prone; yet the
commercial pressure to abandon costly SAS licences and to embrace open-source
Python ecosystems has never been stronger.

This opening chapter establishes the ground on which every subsequent design
decision rests. We begin by situating the work within its academic and
industrial frames, then move to a precise characterisation of the migration
problem, a critical survey of existing tools, and a description of the solution
we chose to build. The chapter closes with a discussion of the methodological
framework that guided the fifteen-week development sprint from first prototype
to production-ready system.

Understanding the \emph{why} before the \emph{how} is not merely a formal
requirement of the PFE report genre; it is epistemically necessary. The
architectural choices examined in Chapters~\ref{ch:architecture} and
\ref{ch:implementation} --- eight orchestrated agents, three \ac{RAG} tiers,
\ac{SMT}-based formal verification, and a four-layer semantic scoring framework
--- each trace back to a concrete deficiency identified in this chapter's
problem statement.

% ---------------------------------------------------------------------------
\section{Academic Institution}
\label{sec:institution}
% ---------------------------------------------------------------------------

This project was carried out in partial fulfilment of the requirements for an
Engineering Degree at \textbf{[Institution Name]}, Faculty of
\textbf{[Faculty / Department]}. The programme combines rigorous theoretical
foundations in software engineering, artificial intelligence, and distributed
systems with a mandatory industry-facing end-of-study project (PFE). The
academic supervisor provided guidance on research methodology, evaluation
design, and the scientific framing of results, while the industrial host
company shaped the problem definition, feasibility constraints, and acceptance
criteria that governed the implementation.

The project was evaluated by a jury comprising a president, an examiner, and
the academic supervisor (see title page). The jury assessment criteria
emphasised not only working software but also the quality of the scientific
contribution: the novelty of the proposed solution relative to the state of
the art, the rigour of the evaluation methodology, and the clarity with which
design trade-offs are argued.

% ---------------------------------------------------------------------------
\section{Host Company}
\label{sec:host-company}
% ---------------------------------------------------------------------------

\subsection{Company Profile and Business Context}
\label{subsec:company-overview}

The internship was conducted at \textbf{[Company Name]}, a
\textbf{[sector, e.g.\ financial services / data engineering consultancy]}
organisation whose data engineering teams rely heavily on \ac{SAS} for
\ac{ETL} processing, statistical modelling, and regulatory reporting. Like
many large enterprises that adopted SAS in the 1990s and 2000s, the organisation
faces a strategic inflection point: SAS licence costs have risen steadily, the
talent pool of experienced SAS developers is contracting, and internal
stakeholders increasingly demand Python-native analytics pipelines compatible
with cloud-based data platforms.

[FIGURE 1.1 --- Caption: ``Organisational chart of the host company and the
position of the internship team within the data engineering division.'' ---
Suggested: simplified org-chart with the data-engineering department
highlighted.]

The company's \ac{SAS} codebase is substantial. It spans multiple business
domains, includes cross-file macro libraries, and contains several hundred
thousand lines of production code accumulated over two decades. This scale
makes purely manual migration impractical: estimates from the project sponsor
placed the manual effort at several person-years, with unacceptable risk of
semantic errors being silently introduced during rewriting.

\subsection{Mission, Services, and Technological Focus}
\label{subsec:mission-vision}

The organisation's data engineering division is responsible for maintaining the
analytical pipelines that feed both internal decision-making systems and
external regulatory submissions. Its primary technology stack as of the start of
this internship included SAS 9.4 (batch and interactive), Oracle SQL, and an
emerging Python/PySpark layer deployed on Azure Databricks. The migration
initiative was therefore not an isolated experiment but a core component of a
multi-year modernisation roadmap.

From a technological standpoint, the host company had already made a strategic
commitment to the Microsoft Azure ecosystem, including Azure Blob Storage for
file persistence, Azure Queue Storage for asynchronous job dispatch, Azure
Container Apps for containerised workloads, and Azure Key Vault for secrets
management. This pre-existing infrastructure commitment directly influenced
several architectural choices described in Chapter~\ref{ch:architecture} ---
notably the cloud-native deployment design of \codara{}.

\subsection{Internship Scope}
\label{subsec:internship-scope}

The internship ran for approximately fifteen weeks, from early January to
mid-April 2026. The agreed scope covered the complete design, implementation,
testing, and documentation of an automated SAS-to-Python conversion accelerator,
with the following boundaries explicitly negotiated at the outset:

\begin{itemize}
  \item \textbf{In scope}: SAS Base (DATA steps, PROC steps, macro language
    constructs, cross-file dependencies), translation to standard Python with
    pandas and NumPy, and partial support for PySpark output where explicitly
    flagged by the user.
  \item \textbf{Out of scope}: SAS/STAT procedures beyond PROC MEANS and
    PROC FREQ, SAS/IML matrix algebra, SAS/GRAPH visualisation procedures,
    and interactive SAS sessions.
  \item \textbf{Quality bar}: a working, deployable system with a documented
    evaluation against a gold-standard corpus, not merely a proof-of-concept
    script.
\end{itemize}

The scope boundaries were not arbitrary. They reflected an 80/20 analysis of
the company's actual codebase: the vast majority of production SAS files fall
within Base SAS, and the tail of exotic procedures was deemed acceptable to
handle as an exception via the system's partial-conversion pathway.

% ---------------------------------------------------------------------------
\section{Project Presentation}
\label{sec:project-presentation}
% ---------------------------------------------------------------------------

\subsection{Initial Situation and Problem Statement}
\label{subsec:problem}

The technical problem can be stated precisely. A SAS program is a sequence of
\emph{steps} --- DATA steps that manipulate datasets row by row and PROC steps
that invoke built-in analytical procedures. The language is inherently
procedural: control flows through implicit program data vectors, RETAIN
statements carry values across observations, and \texttt{FIRST.}/\texttt{LAST.}
variables detect group boundaries. None of these constructs has a direct Python
equivalent; each requires a semantic-preserving transformation that is
non-trivial to automate.

Three layers of difficulty compound the core translation problem:

\begin{enumerate}
  \item \textbf{Macro meta-programming.} SAS macros are a preprocessor language
    that generates SAS code before execution. A macro call such as
    \texttt{\%apply\_etl(input=raw, output=clean, year=2024)} may expand into
    hundreds of lines of DATA-step code whose structure depends entirely on the
    macro definition and its parameter values. Resolving macro dependencies
    across files and correctly expanding their logic is the hardest single
    problem in SAS migration.

  \item \textbf{Cross-file dependencies.} Enterprise SAS programmes are rarely
    self-contained. They \texttt{\%INCLUDE} external macro libraries, reference
    datasets by libref, and chain multiple programmes through agreed intermediate
    dataset names. Any translation tool that treats files in isolation will
    silently break these inter-file contracts.

  \item \textbf{Semantic correctness at scale.} A translator that produces
    syntactically valid Python code is not necessarily correct. Missing-value
    arithmetic (SAS missing versus pandas \texttt{NaN}), BY-group processing
    semantics, and PROC SQL join behaviour all differ subtly between the two
    languages. Without a verification mechanism, semantic errors accumulate
    invisibly in the translated output.
\end{enumerate}

These three dimensions --- macro complexity, cross-file scope, and semantic
correctness --- form the core requirements that any credible solution must
address, and they are directly reflected in the architecture of \codara{}.

[FIGURE 1.2 --- Caption: ``Three-layer complexity model for SAS migration:
structural (macro meta-programming), topological (cross-file dependencies),
and semantic (behavioural equivalence).'' --- Suggested: three nested rings
or a three-column table with example constructs in each layer.]

\subsection{Study and Critique of Existing Solutions}
\label{subsec:existing-solutions}

Before committing to a bespoke solution, we surveyed the landscape of available
tools along two axes: commercial migration products and academic
research approaches.

\textbf{Commercial tools.} Several vendors offer SAS-to-Python migration
assistants. Products such as \textit{Analytium SAS Converter},
\textit{WhereScape} migration utilities, and bespoke consulting engagements
typically rely on pattern-matching rule sets. While effective for simple DATA
steps and common PROCs, rule-based systems fail on macro-heavy code --- they
cannot expand parameterised macros without execution --- and they produce no
confidence signal alongside the translation. The user has no way to know, from
the tool's output alone, whether the translated code is semantically equivalent
to the original. Furthermore, these tools are inherently single-file: they
process each SAS programme in isolation without resolving cross-file macro
dependencies.

\textbf{\ac{LLM}-based approaches.} The emergence of code-capable \acp{LLM}
such as GPT-4 and Code LLaMA opened a new avenue: asking the model to translate
SAS snippets directly. Early experiments within the host company confirmed that
raw LLM translation is surprisingly effective for isolated, low-complexity DATA
steps but degrades rapidly for longer files, macro-parameterised constructs, and
cross-file references. The model has no persistent memory across translation
calls, no mechanism to verify its output, and no way to retrieve relevant
company-specific translation patterns from prior work. Token-limit constraints
also prevent whole-file translation of large SAS programmes.

\textbf{Research tools.} Academic literature on automated code migration has
produced systems such as \textit{CodeTransOcean} and various
transpiler-based approaches~[1, 2], but these are generally evaluated on
small, self-contained code snippets and do not address the enterprise-scale
challenges of macro resolution, cross-file dependencies, or formal semantic
verification.

[TABLE 1.1 --- Caption: ``Comparison of existing SAS migration approaches.''
--- Columns: Tool/Approach, Macro Resolution, Cross-File Support, Semantic
Verification, Output Confidence Signal, Open Source.]

The critical gap that this analysis revealed is not merely a matter of
translation accuracy. It is the absence of any tool that simultaneously
addresses all three dimensions of the problem: macro-aware parsing, cross-file
dependency resolution, and \emph{verifiable} semantic correctness. \codara{}
was designed to fill precisely this gap.

\subsection{Proposed Solution Overview}
\label{subsec:proposed-solution}

\codara{} is a multi-agent \ac{SAS}-to-Python/PySpark conversion accelerator.
It takes one or more SAS source files as input and produces equivalent Python
scripts accompanied by a structured validation report. Internally, the system
routes each input through an eight-node \ac{LangGraph} pipeline in which
specialised agents handle distinct concerns in sequence: file parsing and
cross-file dependency resolution, token-level streaming and boundary detection,
\ac{RAPTOR} semantic clustering, complexity scoring and strategy routing,
persistence and dependency indexing, three-tier \ac{RAG}-augmented translation,
and finally script merging and report generation.

The system version delivered at the close of this internship is \textbf{v3.1.0},
which adds Z3 \ac{SMT} formal verification, the \ac{CDAIS} adversarial test
synthesis framework, and the \textit{SemantiCheck} four-layer semantic scoring
system on top of the core translation pipeline. The knowledge base comprises
330 verified SAS–Python example pairs stored in a 768-dimensional vector index.
The full test suite contains 309 unit and integration tests (0 failures in the
Week 15 CI run), and the system received an audit grade of A- following four
rounds of security and code-quality review.

\codara{} is not a point solution. It exposes a documented REST API, a
React-based web dashboard, a Docker Compose deployment configuration, and a
six-job GitHub Actions \ac{CI/CD} pipeline targeting Azure Container Apps.
This production orientation was a deliberate project requirement, not an
afterthought.

[FIGURE 1.3 --- Caption: ``High-level overview of the \codara{} system: a
multi-agent LangGraph pipeline bridging SAS source files and validated
Python/PySpark output.'' --- Suggested: block-flow diagram from SAS input
through the 8-node pipeline to Python output + validation report.]

\subsection{Expected Deliverables}
\label{subsec:deliverables}

The deliverables negotiated with the host company and the academic supervisor
at project kick-off were as follows:

\begin{enumerate}
  \item A working backend service implementing the eight-node translation
    pipeline, exposed via a documented \ac{REST} \ac{API}.
  \item A React web frontend providing upload, conversion monitoring,
    side-by-side diff view, and admin panels for knowledge-base management,
    audit logs, and system health.
  \item A gold-standard evaluation corpus of 50 annotated SAS files
    (45 with hand-verified Python translations), covering 721 individual blocks
    across three difficulty tiers (basic, medium, hard), used to measure
    boundary detection accuracy, retrieval quality, and translation success
    rate.
  \item A formal evaluation report quantifying performance against the targets
    established in the project specification (see Chapter~\ref{ch:specifications}).
  \item Production deployment artefacts: Dockerfile, Docker Compose
    configuration, GitHub Actions pipeline, and Azure Key Vault integration.
  \item This end-of-study report documenting all design decisions, theoretical
    foundations, experimental results, and critical assessment.
\end{enumerate}

All six deliverables were completed within the fifteen-week sprint. The status
of each is detailed in Chapter~\ref{ch:evaluation}.

% ---------------------------------------------------------------------------
\section{Methodology Orientation}
\label{sec:methodology}
% ---------------------------------------------------------------------------

\subsection{Overview of Candidate Methodologies}
\label{subsec:methodologies-overview}

Choosing a methodology for a project of this kind requires balancing two
competing demands. On one hand, the project is a research artefact: it
proposes novel solutions (three-tier \ac{RAG}, \ac{CDAIS}, SemantiCheck) that
must be situated within and compared against a body of prior work. On the other
hand, it is a software engineering deliverable with a hard completion deadline
and an industrial client who needs a working system. Pure research methodologies
such as \ac{CRISP-DM} or the Hypothetico-Deductive model are ill-suited to
iterative software construction; pure agile methods such as Scrum provide no
scaffolding for scientific evaluation and knowledge contribution.

Three candidate methodologies were considered:

\begin{itemize}
  \item \textbf{Scrum / Agile}: strong on iterative delivery and customer
    feedback, weak on scientific rigour and reproducibility of experiments.
  \item \textbf{\ac{CRISP-DM}}: appropriate for data-mining projects with
    well-defined datasets and models, but not designed for the construction
    of novel software artefacts.
  \item \textbf{\ac{DSR}}: a research paradigm from the Information Systems
    discipline~[3] that frames the research product as a designed artefact and
    requires its evaluation against explicit utility claims. \ac{DSR} explicitly
    accommodates the construction-plus-evaluation cycle needed here.
\end{itemize}

\subsection{Design Science Research (DSR)}
\label{subsec:dsr}

\ac{DSR}, as codified by Hevner et al.~[3], provides a seven-guideline
framework for research that produces and evaluates designed artefacts. Its
central tenet is that the artefact must be shown to solve a \emph{relevant
problem} better than existing alternatives, and this demonstration must meet
\emph{research rigour} standards --- reproducible experiments, explicit metrics,
and honest critical assessment of limitations.

In the \ac{DSR} framing, \codara{} is a \emph{system artefact}: a fully
implemented software system whose utility is evaluated against measurable
performance targets. The five design cycles identified in \ac{DSR} map
naturally onto the project's week-by-week sprints:

\begin{enumerate}
  \item \textbf{Problem awareness}: the migration problem and its three-layer
    complexity characterised in Section~\ref{subsec:problem}.
  \item \textbf{Suggestion}: the multi-agent, \ac{RAG}-augmented, verification-
    backed architecture proposed as the solution.
  \item \textbf{Development}: fifteen weeks of iterative implementation
    (Chapters~\ref{ch:architecture} and~\ref{ch:implementation}).
  \item \textbf{Evaluation}: gold-standard benchmarking, ablation studies, and
    torture-test execution (Chapter~\ref{ch:evaluation}).
  \item \textbf{Conclusion}: generalisation of findings and future-work
    directions (Chapter~\ref{ch:evaluation}, Section~\ref{sec:future-work}).
\end{enumerate}

\ac{DSR} also mandates that novelty claims be explicitly stated and argued.
Section~\ref{subsec:novelty} in Chapter~\ref{ch:evaluation} enumerates the
five distinct novelty claims of this work and situates each relative to prior
literature.

\subsection{Adopted Hybrid Methodology: DSR + Agile}
\label{subsec:hybrid-methodology}

In practice, a pure \ac{DSR} approach would front-load the full system design
before any implementation begins --- a luxury incompatible with a fifteen-week
deadline and the inherent uncertainty of working with \acp{LLM} whose
capabilities are only revealed through empirical probing. We therefore adopted
a hybrid methodology that uses \ac{DSR} for the research framing and Agile
(specifically two-week sprint cycles) for the development execution.

[FIGURE 1.4 --- Caption: ``Hybrid DSR + Agile methodology: outer DSR cycle
governing problem-solution-evaluation, inner Agile sprints driving
iterative implementation.'' --- Suggested: concentric rings or two-lane
swimlane diagram with DSR on the outer ring and numbered Agile sprints on
the inner.]

The sprint structure was as follows. Weeks 1--2 focused on file analysis
and cross-file dependency resolution (pipeline layer L2-A). Weeks 3--4
delivered the FSM-based streaming parser and boundary detector (L2-B and
L2-C). Weeks 5--6 added RAPTOR clustering and complexity scoring (L2-C/D).
Weeks 7--8 implemented orchestration, persistence, and the audit infrastructure
(L2-E and orchestration). Week 9 introduced the rate-limiting, circuit-breaker,
and knowledge-base subsystems. Weeks 10--11 completed the three-tier
translation and merge pipeline (L3 and L4). Week 12 built the ablation
evaluation infrastructure. Week 13 performed a major architectural consolidation
(11 nodes reduced to 8) and introduced CI/CD, Docker, and CodeQL security
scanning. Weeks 14--15 added Z3 formal verification, \ac{CDAIS}, and
SemantiCheck, and produced the final evaluation dataset.

The \ac{DSR} evaluation cycle was not a one-off end-of-project exercise. Each
sprint produced a measurable artefact --- a benchmark run, an ablation result,
a test coverage report --- that fed back into the design of the next sprint.
This tight coupling between development and evaluation is precisely what
distinguishes the hybrid approach from naive iterative development, and it is
what allowed the project to claim not only a working system but a rigorously
evaluated one.

[TABLE 1.2 --- Caption: ``Sprint schedule and key deliverables.'' --- Columns:
Weeks, Layer, Key Deliverable, Cumulative Tests.]

% ---------------------------------------------------------------------------
\section{Conclusion}
\label{sec:ch1-conclusion}
% ---------------------------------------------------------------------------

This chapter has laid the contextual foundation for everything that follows.
We identified a genuine, industry-scale problem --- the automated migration of
enterprise SAS code to Python --- characterised its three-layer technical
complexity (macro meta-programming, cross-file topology, and semantic
equivalence), and demonstrated through a comparative survey that no existing
tool addresses all three layers simultaneously. The proposed solution,
\codara{}, is a multi-agent, \ac{RAG}-augmented, formally verified conversion
accelerator delivered as a production-ready system over a fifteen-week sprint.

The methodological choice of a hybrid \ac{DSR}+Agile approach was deliberate:
it provides the scientific scaffolding needed to frame, evaluate, and situate
the novelty of the work, while preserving the iterative flexibility required
to navigate the empirical uncertainties inherent in \ac{LLM}-based systems.

The theoretical foundations that underpin the architectural choices are examined
in the following chapter.

% ---------------------------------------------------------------------------
% References for Chapter 1
% (IEEE numeric, sorted by order of appearance)
% ---------------------------------------------------------------------------
%
% [1] Agashe, K. et al., "SWE-bench: Can Language Models Resolve Real-World
%     GitHub Issues?", ICLR 2024.
% [2] Zhu, M. et al., "CodeTransOcean: A Comprehensive Multilingual Benchmark
%     for Code Translation", EMNLP 2023 Findings.
% [3] Hevner, A. R. et al., "Design Science in Information Systems Research",
%     MIS Quarterly, vol. 28, no. 1, pp. 75-105, Mar. 2004.
